\documentclass[11pt]{article}
\usepackage[T1]{fontenc}
\usepackage{amsmath,amssymb,amsthm}
\usepackage{fourier}
\usepackage{booktabs}
\usepackage{enumitem}
\usepackage[margin=1in]{geometry}
\usepackage[hidelinks]{hyperref}

\newcommand{\code}[1]{\texttt{#1}}
\setlength{\emergencystretch}{3em}

\title{Minimal sets for the classes modulo 2, 3, 4, 6 and 8\\ in bases with at most two terminal digits}
\author{}  % add author
\date{19 November 2026, 23:52 MST (UTC$-$07:00)\\ Updated 20 November 2026, 01:40 MST}

\begin{document}
\maketitle

\begin{abstract}
For the fifteen bases $b$ with at most two possible terminal digits ($c(b)\le 2$; see \code{terminal-digits.tex}) and the eleven classes $a \bmod m$ with $m\in\{2,3,4,6,8\}$ and $\gcd(a,m)=1$, we attempted all $165$ minimal sets. $96$ are complete, verified independently, and stored in \code{data-computed/}. $20$ are reduced to simple families that remain open after searching to length $10{,}000$. In the other $49$, \code{minimal} did not finish within the time available. Section~\ref{sec:update} describes the changes since the first version of this report. The report records the status of each set, the open families, and what the unfinished cases need.
\end{abstract}

\section{Scope}

The bases are those with $c(b)\le 2$:
\[
  b\in\{2,3,4,6,8,12,18,24,30\}\ (c=1),\qquad b\in\{5,10,14,20,42,60\}\ (c=2).
\]
The classes are $1 \bmod 2$; $1,2 \bmod 3$; $1,3 \bmod 4$; $1,5 \bmod 6$; and $1,3,5,7 \bmod 8$. For each pair the target is the set of minimal elements, under the subsequence order, of the base-$b$ representations of the primes in the class (\code{bright-algorithm.tex}, Section~10).

\section{Status}

\begin{table}[h]
\centering\small
\setlength{\tabcolsep}{4.5pt}
\begin{tabular}{@{}r|r|rr|rr|rr|rrrr@{}}
\toprule
 & \multicolumn{1}{c|}{mod 2} & \multicolumn{2}{c|}{mod 3} & \multicolumn{2}{c|}{mod 4} & \multicolumn{2}{c|}{mod 6} & \multicolumn{4}{c}{mod 8}\\
$b$ & 1 & 1 & 2 & 1 & 3 & 1 & 5 & 1 & 3 & 5 & 7\\
\midrule
2 & 1 & 1 & 1 & 1 & 1 & 1 & 1 & 2 & 1 & 1 & 1 \\
3 & 4 & 2 & 1 & 5 & 2 & 2 & 1 & 29 & 4 & 3 & 8 \\
4 & 3 & $\circ_{1}$ & 2 & 3 & 1 & $\circ_{1}$ & 3 & 9 & 1 & 2 & 2 \\
5 & 9 & 49 & 7 & 80 & 4 & 49 & 19 & $\circ_{14}$ & 60 & $\circ_{6}$ & $\circ_{19}$ \\
6 & 7 & 7 & 2 & 10 & 6 & 7 & 1 & 75 & 8 & 7 & 10 \\
8 & 15 & $\circ_{5}$ & $\times$ & 38 & 2 & $\circ_{5}$ & $\times$ & 38 & 1 & 1 & 1 \\
10 & 33 & $\circ_{2}$ & 251 & 146 & $\circ_{1}$ & $\circ_{2}$ & 309 & $\circ_{3}$ & $\circ_{2}$ & $\circ_{1}$ & $\circ_{2}$ \\
12 & 22 & 40 & 3 & 52 & 3 & 40 & 2 & $\circ_{1}$ & 132 & 63 & 103 \\
14 & 241 & $\times$ & $\times$ & $\circ_{1}$ & $\circ_{1}$ & $\times$ & $\times$ & $\times$ & $\times$ & $\times$ & $\times$ \\
18 & 50 & 111 & 4 & $\circ_{2}$ & $\circ_{3}$ & 111 & 3 & $\times$ & $\times$ & $\times$ & $\times$ \\
20 & 917 & $\times$ & $\times$ & 1{,}380 & 4 & $\times$ & $\times$ & $\times$ & $\circ_{2}$ & $\times$ & $\times$ \\
24 & 360 & 811 & 5 & 648 & 5 & 811 & 4 & 891 & 3 & 2 & 2 \\
30 & 220 & 489 & 6 & $\times$ & $\times$ & 489 & 5 & $\times$ & $\times$ & $\times$ & $\times$ \\
42 & $\times$ & $\times$ & $\times$ & $\times$ & $\times$ & $\times$ & $\times$ & $\times$ & $\times$ & $\times$ & $\times$ \\
60 & $\times$ & $\times$ & $\times$ & $\times$ & $\times$ & $\times$ & $\times$ & $\times$ & $\times$ & $\times$ & $\times$ \\
\bottomrule
\end{tabular}
\caption{Status of the 165 sets. A number is the size of a complete set; $\circ_k$ means $k$ simple families are still open after searching to length $10{,}000$; $\times$ means \code{minimal} did not finish.}
\end{table}

\begin{center}
\begin{tabular}{@{}lrp{8.2cm}@{}}
\toprule
Status & Sets & Where\\
\midrule
Complete & 96 & all classes in bases 2, 3, 6, 24; most in 4, 5, 12; some in 8, 10, 14, 18, 20, 30\\
Open simple families & 20 & see Table~\ref{tab:open}\\
\code{minimal} unfinished & 49 & bases 42 and 60 entirely; the rest in 8, 14, 18, 20, 30\\
\bottomrule
\end{tabular}
\end{center}

\section{Method}

Each set was computed in its own folder with \code{minimal B --class a m} (time limit 20\,s, then 150\,s for sets in bases below 42), \code{simple 1 300} and \code{organize}. Where families remained, this was followed by \code{sieve-start 300 3000}, \code{sieve --pmax 1e6} and \code{search 300 3000} (time limit 240\,s). The code was \code{mepn-python} at commit \code{5f7a5a5}, with the built-in Baillie--PSW test and srsieve 0.6.17, on two processor cores. The second round (Section~\ref{sec:update}) used commit \code{ae2a7b7}, searched to length $10{,}000$ with \code{sieve --pmax 1e7}, and tested primality with GMP. The full table of complete sets, with the lengths searched, is in \code{data-computed/README.md}.

\section{Verification}

Each of the 92 complete sets was checked by a separate script:
\begin{enumerate}[label=(\arabic*)]
\item every entry is a probable prime in the class;
\item no entry is a subsequence of another;
\item every prime $p<2\cdot10^7$ in the class contains an entry;
\item every element of the published minimal set for base $b$ (in \code{data-minimal-sets/}) that lies in the class appears in the computed set.
\end{enumerate}
Check (4) uses independent data. If $p$ is minimal among all primes and $p\equiv a\pmod m$, then $p$ is also minimal among the primes $\equiv a \pmod m$, so it must be in the computed set. Published sets exist for every base here except 60, so all 92 complete sets were checked this way. All four checks passed for every set.

\section{Relations between classes}

\begin{itemize}
\item \textbf{$1 \bmod 6 = 1 \bmod 3$.} Every prime $\equiv 1 \pmod 3$ is odd, so the two classes contain the same primes. The computed sets agree in every base where both finished.
\item \textbf{$5 \bmod 6$ versus $2 \bmod 3$.} The primes $\equiv 5 \pmod 6$ are those $\equiv 2 \pmod 3$ except $2$. The $5 \bmod 6$ set is the $2 \bmod 3$ set with the word for $2$ removed and the primes that contained it as their only minimal subword added; in base 10 these include 2087 and 20021.
\item \textbf{$1 \bmod 2$.} The same reasoning gives the odd primes: the published set without $2$, plus the odd primes that only $2$ ruled out. This was confirmed in all 13 bases where the set finished, with between 0 (base 2) and 267 (base 20) added elements, each containing the word for $2$.
\end{itemize}

\section{Open simple families}

\begin{table}[h]
\centering\small
\begin{tabular}{@{}rlrp{6.3cm}r@{}}
\toprule
$b$ & class & primes & open families & searched to\\
\midrule
4 & 1 mod 3 (= 1 mod 6) & 17 & \code{1*21} & 10{,}000\\
5 & 1 mod 8 & 776 & 14 families, e.g.\ \code{10*11123}, \code{1*33} & 10{,}000\\
5 & 5 mod 8 & 343 & \code{134*1}, \code{343*}, \code{34443*}, \code{43*}, \code{4*113}, \code{4*324} & 10{,}000\\
5 & 7 mod 8 & 320 & 19 families, e.g.\ \code{230*33}, \code{3*1} & 10{,}000\\
8 & 1 mod 3 (= 1 mod 6) & 259 & \code{212*1}, \code{2*1}, \code{2*51}, \code{40*11}, \code{40*2211} & 10{,}000\\
10 & 1 mod 3 (= 1 mod 6) & 1{,}028 & \code{8*69}, \code{8*9} & 10{,}000\\
10 & 3 mod 4 & 112 & \code{2*99} & 10{,}000\\
10 & 1 mod 8 & 905 & \code{10*1}, \code{6*2321}, \code{7*5049} & 7{,}977\\
10 & 3 mod 8 & 438 & \code{2*99}, \code{5*451} & 10{,}000\\
10 & 5 mod 8 & 868 & \code{1*221} & 10{,}000\\
10 & 7 mod 8 & 1{,}015 & \code{2*519}, \code{80*8511} & 8{,}710\\
12 & 1 mod 8 & 182 & \code{10*1} & 10{,}000\\
14 & 1 mod 4 & 1{,}401 & \code{A*6BB} & 7{,}699\\
14 & 3 mod 4 & 1{,}243 & \code{4D*} & 7{,}828\\
18 & 1 mod 4 & 1{,}178 & \code{10*1}, \code{G*077} & 7{,}855\\
18 & 3 mod 4 & 2{,}289 & \code{40*8HD}, \code{C*D1}, \code{E50*F5} & 6{,}698\\
20 & 3 mod 8 & 2{,}282 & \code{1*817}, \code{D*GD7} & 5{,}882\\
\bottomrule
\end{tabular}
\caption{The 20 sets with open families (the mod 6 duplicates share a row). ``Primes'' counts the minimal primes found so far; ``searched to'' is the largest repeat count $n$ tested. Any further minimal prime in such a family has more than that many repeated digits.}
\label{tab:open}
\end{table}

Three of the open families are \code{10*1}, in bases 10 ($1 \bmod 8$), 12 ($1 \bmod 8$) and 18 ($1 \bmod 4$). Its members are $1\,0^k\,1=b^{k+1}+1$, which can be prime only when $k+1$ is a power of $2$, so these searches reduce to generalized Fermat numbers $b^{2^j}+1$. In base 10 this is exactly the family discussed in \code{terminal-digits.tex}: the known primes $11$ and $101$ are $3$ and $5 \bmod 8$, while $10^k+1\equiv 1 \pmod 8$ for $k\ge 3$. So the set for $1 \bmod 8$ in base 10 is complete if and only if one decides whether some $10^{2^j}+1$ with $j\ge 2$ is prime, which is an open problem. The same may hold in bases 12 and 18.

The other families need longer searches. With GMP a test takes about 2.4\,s at 10{,}000 digits and 12\,s at 20{,}000, so searches to about 20{,}000 are practical in the environment used. Beyond that, LLR is the better tool for these forms; \code{simple - M} and \code{search -} continue across sessions.

\section{Unfinished runs of \code{minimal}}

The sets where \code{minimal} did not finish behaved in three distinct ways. The third kind has since been resolved (Section~\ref{sec:update}).

\begin{itemize}
\item \textbf{Bases 42 and 60 (22 sets).} \code{minimal} had not finished even its initial stage after 20\,s. For a base of this size, the initial kernel (all minimal words of up to three digits) and the $b(b-1)$ starting families are already large, and every later step is correspondingly expensive. These need a long run on a dedicated machine.
\item \textbf{Growth (27 sets):} base 8 at $2 \bmod 3$ and $5 \bmod 6$; base 14 at all classes mod 3, 6 and 8; base 18 at all classes mod 8; base 20 at all classes mod 3 and 6 and at $1,5,7 \bmod 8$; base 30 at all classes mod 4 and 8. After 150\,s, these had between about 1{,}600 and 20{,}000 open families after only 0--12 iterations, and the four mod-8 classes of base 30 had not yet finished the initial stage. A sparser target set means fewer kernel words, so fewer digits are ruled out and exploration multiplies the families. These may finish with hours of computation, resuming with \code{minimal resume}.
\item \textbf{Circling (6 sets, now resolved):} base 5 with $1,2 \bmod 3$ and $1,3,5,7 \bmod 8$. After 150\,s, these had reached iterations 77--296 with only 600--4{,}300 families. The family count was not exploding but the loop did not end: the non-termination described in \code{bright-algorithm.tex}, Section~9, caused by non-simple families that no certificate could rule out. Section~\ref{sec:update} identifies them and gives the certificate that now does.
\end{itemize}

\section{Update, 20 November 2026}\label{sec:update}

\paragraph{The circling base-5 sets.} The six base-5 sets were each stuck on one or two non-simple families that kept regenerating, such as $11\,0^*4^*1$ for $3 \bmod 8$. Its members are $N=6\cdot5^{a+b+1}+5^{b+1}-4$, with $a$ zeros and $b$ fours. If $b$ is odd, then $5^{b+1}\equiv 1$ and $N\equiv 0 \pmod 3$. If $b$ is even, then $N\equiv 7 \pmod 8$, outside the class. So the family contains no prime $\equiv 3 \pmod 8$. But the program tested ``outside the class'' and ``divisible by 2, 3 or 5'' separately, and each half of the family passes a different one of the two tests.

The class test now computes the reachable residues modulo $\operatorname{lcm}(m,30)$ and discards a family when no residue is both $\equiv a \pmod m$ and coprime to $30$ (\code{bright-algorithm.tex}, Section~5.6). With this, all six base-5 sets finish \code{minimal} in under 10\,s. Three are complete ($1, 2 \bmod 3$ and $3 \bmod 8$). The other three are down to simple families, which remain open after searching to length $10{,}000$. Base 10, $1 \bmod 8$, which had been growing, now also finishes \code{minimal}.

\paragraph{Regression check.} All 92 sets that were complete before were recomputed with the new code, and every one came out identical. Unrestricted runs do not use the class test and are unaffected.

\paragraph{Primality testing with GMP.} gmpy2 cannot be installed in the environment used, but the GMP library is present. The program now calls \code{mpz\_probab\_prime\_p(n, 25)} in it directly, which is the routine used by gmpy2 and by \code{minimal.c}. A test on a 10{,}000-digit number takes about 2.4\,s instead of about 40\,s. With it, the open searches were extended from length 3000 to 10{,}000. Base 12, $3 \bmod 8$ was completed: its last family, \code{8*77}, gave the prime $8^{7733}77$ (7{,}735 base-12 digits, 8{,}348 decimal digits), confirmed by both GMP and the built-in Baillie--PSW test.

\paragraph{New complete sets.} Base 5 at $1 \bmod 3$ (49 primes), $2 \bmod 3$ (7) and $3 \bmod 8$ (60), and base 12 at $3 \bmod 8$ (132). All four passed the checks of Section~4.

\section{Next steps}

\begin{enumerate}[label=\arabic*.]
\item Continue the 20 open searches beyond length $10{,}000$, here with GMP and on a machine with LLR.
\item Run the 27 growing sets with \code{minimal resume} over several hours, and bases 42 and 60 as long-running jobs.
\end{enumerate}
New complete sets go into \code{data-computed/} with the same checks.

\end{document}
